\section{نتیجه‌گیری}
\label{sec:conclusion}

این پروژه یک معماری کامل برای بارگذاری و کوئری‌گرفتن از مجموعه‌داده‌ی \lr{IMDb} روی \lr{PostgreSQL 16} ساخت: یک صف پیام با \lr{SELECT ... FOR UPDATE SKIP LOCKED} برای بارگذاری هم‌زمان و بدون رقابت، بسته‌بندی دسته‌ای برای رد کردن آستانه‌ی فشرده‌سازی \lr{TOAST}، هشت سناریوی کوئری طبق سند پروژه، و شش ایندکس که هرکدام به یک شرط مشخص در این سناریوها وصل است.

هر بخش از این پروژه -- اسکیمای پایگاه‌داده، توابع صف، تولیدکننده و مصرف‌کننده، کوئری‌ها، و طراحی ایندکس -- حداقل یک‌بار با یک بازبینی چندعاملی بررسی شده و باگ‌های واقعی پیداشده در آن اصلاح شده‌اند. برای بخش‌هایی که نیازی به اجرای واقعی نداشتند (نحو \lr{SQL}، منطق \lr{Python}، تست‌های خودکار)، درستی با ابزارهایی مثل \lr{py\_compile} و آزمون دستی چک شده است.

\subsection{محدودیت این نگارش}

همه‌ی عددهای بخش \ref{sec:indexing} از یک اجرای واقعی \lr{PostgreSQL 16} داخل \lr{Docker} به دست آمده‌اند، نه از پیش‌بینی. تنها محدودیت واقعی این اجرا فضای دیسک بود: لپ‌تاپی که این پروژه رویش اجرا شد فقط حدود \lr{30} گیگابایت فضای خالی داشت، و بیشتر از این هم نمی‌شد جا باز کرد. دیتاست کامل \lr{IMDb} -- مخصوصاً جدول \lr{title\_principals} با نزدیک \lr{90} میلیون سطر -- توی این فضا جا نمی‌شد. برای همین فقط همین یک جدول با فلگ \lr{--limit-rows} به \lr{8} میلیون سطر (حدود \lr{9\%}) محدود شد؛ شش جدول دیگر کامل بارگذاری شدند (جزئیات در بخش \ref{sec:real-data-scale}). نتیجه‌ی عملی این محدودیت این است که سناریوی \lr{5} -- تنها سناریویی که از \lr{title\_principals} می‌خواند -- روی حجمی کوچک‌تر از واقعیت سنجیده شده؛ بهبود واقعی آن روی دیتاست کامل احتمالاً بیشتر از عدد گزارش‌شده است.

یک نتیجه‌ی دیگر هم ارزش گفتن دارد: ایندکس \lr{GIN} ترکیبی، سناریوی \lr{1} را کندتر کرد نه سریع‌تر. این عدد را حذف یا نرم نکردیم، چون خودش یک درس واقعی از این پروژه است. ایندکس‌گذاری فقط با نگاه‌کردن به لیست شرط‌های کوئری کافی نیست. باید با \lr{EXPLAIN (ANALYZE, BUFFERS)} واقعی سنجید، چون گزینش‌پذیری کم یک شرط -- مثل \lr{title\_type = 'movie'} -- می‌تواند نتیجه‌ی برعکس بدهد.

\subsection{کارهای آینده}

اگر بخواهیم این پروژه را گسترش دهیم، سه مسیر طبیعی برای ادامه وجود دارد: افراز کردن (\lr{partitioning}) جدول \lr{title\_principals} بر اساس \lr{category} برای کم‌کردن بیشتر حجم اسکن‌شده در سناریوی \lr{5}؛ استفاده از دیدهای مادی‌شده (\lr{materialized views}) برای سناریوهایی که به‌طور دوره‌ای اجرا می‌شوند نه بلادرنگ؛ و پایش مستمر عملکرد با \lr{pg\_stat\_statements} برای پیدا کردن خودکار کوئری‌هایی که با تغییر داده در طول زمان، به ایندکس یا تنظیم جدیدی نیاز پیدا می‌کنند.
